\chapter{Conclusion}

In the title of this thesis we ask "Safe for whom?", reflecting the concern that efforts to protect users from risks when interacting with AIs may only hold up for certain demographics. We coined this concern "safety discrimination" and differentiated between \textit{methods} to measure and provide safeguards being developed without demographic diversity in mind and the actual \textit{behaviour} of the model being discriminatory by failing to protect all users equally. To investigate this concern we mapped the demographic diversity in 23 evaluations and benchmarks concerned with user safety and found this diversity to be mainly absent: only two of the 23 benchmarks systematically cover more than three demographic factors, and only three stand out for meaningful demographic or linguistic breadth at all. To alleviate this observed methodological safety discrimination in safety evaluations, we continued with four different case studies exploring methods to enhance the linguistic and dialectal diversity, and demographic contextualisation. Finally, we also provided a case study of cross-lingual generalisation of linear probes for judge-free detection of social sycophancy. 

In the attempt to answer our question "Safe for whom?", we conclude that we cannot yet say. This inability to draw a verdict, because the tools to answer it do not exist, is itself evidence for the methodological safety discrimination we set out to map. Future work is needed to refine the methods sketched here and to apply them at scale and close to the frontier so that eventually the answer to "Safe for whom?" can be "safe for everyone", with no user left unprotected because of their demographics.

